\documentclass[12pt,a4paper]{article}

% ── Packages ──────────────────────────────────────────────────────────────────
\usepackage[T1]{fontenc}
\usepackage{lmodern}
\usepackage[margin=2.5cm]{geometry}
\usepackage{amsmath,amssymb,amsthm}
\usepackage{mathtools}
\usepackage{bm}
\usepackage{booktabs}
\usepackage{array}
\usepackage{tabularx}
\usepackage{multirow}
\usepackage{xcolor}
\usepackage{graphicx}
\usepackage{float}
\usepackage{tikz}
\usetikzlibrary{arrows.meta, positioning, fit, shapes.geometric, backgrounds, calc}
\usepackage{algorithm}
\usepackage{algpseudocode}
\usepackage{listings}
\usepackage[many]{tcolorbox}
\usepackage{mdframed}
\usepackage{enumitem}
\usepackage{microtype}
\usepackage{setspace}
\usepackage{titlesec}
\usepackage{fancyhdr}
\usepackage{abstract}
\usepackage{caption}
\usepackage{subcaption}

% ── CJK support ───────────────────────────────────────────────────────────────
\usepackage{xeCJK}
\setCJKmainfont{SimSun}[BoldFont=SimHei,ItalicFont=KaiTi]
\setCJKsansfont{SimHei}
\setCJKmonofont{FangSong}

% ── hyperref must be loaded after xeCJK ───────────────────────────────────────
\usepackage{hyperref}
\usepackage{cleveref}

% ── Colors ────────────────────────────────────────────────────────────────────
\definecolor{myblue}{RGB}{31, 78, 121}
\definecolor{myorange}{RGB}{192, 80, 77}
\definecolor{mygreen}{RGB}{0, 112, 0}
\definecolor{mygray}{RGB}{230, 230, 230}
\definecolor{codeteal}{RGB}{0, 128, 128}
\definecolor{codebg}{RGB}{248, 248, 248}

% ── Section formatting ─────────────────────────────────────────────────────────
\titleformat{\section}[block]
  {\large\bfseries\color{myblue}}
  {\thesection.}{0.5em}{}[\vspace{-0.3em}\rule{\linewidth}{0.6pt}]
\titleformat{\subsection}[runin]
  {\normalsize\bfseries\color{myblue!80!black}}
  {\thesubsection.}{0.4em}{}[.\enspace]
\titleformat{\subsubsection}[runin]
  {\normalsize\itshape\color{myblue!70!black}}
  {\thesubsubsection.}{0.4em}{}[.\enspace]

\titlespacing*{\section}{0pt}{1.4ex plus 0.5ex}{0.8ex}
\titlespacing*{\subsection}{0pt}{1ex plus 0.3ex}{0.4ex}

% ── Theorem environments ───────────────────────────────────────────────────────
\newtheorem{definition}{定义}[section]
\newtheorem{principle}{原则}[section]
\newtheorem{observation}{观察}[section]

% ── tcolorbox styles ──────────────────────────────────────────────────────────
\tcbset{
  insight/.style={
    enhanced, breakable,
    colback=myblue!5, colframe=myblue!50!black,
    fonttitle=\bfseries\small, title={核心洞见},
    left=4pt, right=4pt, top=3pt, bottom=3pt,
    arc=2pt
  },
  challenge/.style={
    enhanced, breakable,
    colback=myorange!5, colframe=myorange!70!black,
    fonttitle=\bfseries\small,
    left=4pt, right=4pt, top=3pt, bottom=3pt,
    arc=2pt
  },
  example/.style={
    enhanced, breakable,
    colback=mygreen!5, colframe=mygreen!60!black,
    fonttitle=\bfseries\small,
    left=4pt, right=4pt, top=3pt, bottom=3pt,
    arc=2pt
  }
}

% ── Code listing ──────────────────────────────────────────────────────────────
\lstdefinestyle{dsl}{
  backgroundcolor=\color{codebg},
  basicstyle=\footnotesize\ttfamily,
  keywordstyle=\color{myblue}\bfseries,
  commentstyle=\color{gray}\itshape,
  stringstyle=\color{myorange},
  breaklines=true,
  frame=single,
  framesep=3pt,
  rulecolor=\color{gray!40},
  xleftmargin=8pt,
  xrightmargin=4pt,
  numbers=none,
  captionpos=b,
  morekeywords={sequence, guide, nest, iterate, refine, parallel_merge,
                fold_layers, approximate_in_family, match_factorization,
                coordinate_descent, distribute_divide, distribute_multiply,
                project_to_family, system_to_factor_graph, find_spanning_tree}
}

% ── Fix fancyhdr headheight ──────────────────────────────────────────────────
\setlength{\headheight}{14pt}
\addtolength{\topmargin}{-2pt}

% ── Header/Footer ─────────────────────────────────────────────────────────────
\pagestyle{fancy}
\fancyhf{}
\fancyhead[L]{\small\textit{AlphaDetect：基于自动化算法发现的可解释MIMO检测}}
\fancyhead[R]{\small\thepage}
\fancyfoot[C]{\small\color{gray}研究提案 --- 2026年4月}
\renewcommand{\headrulewidth}{0.4pt}
\renewcommand{\footrulewidth}{0.2pt}

% ── Hyperref setup ────────────────────────────────────────────────────────────
\hypersetup{
  colorlinks=true,
  linkcolor=myblue,
  citecolor=mygreen!70!black,
  urlcolor=codeteal,
  pdftitle={AlphaDetect 研究提案},
  pdfauthor={},
  pdfsubject={MIMO检测, 算法发现, 神经符号AI}
}

% ── Math macros ───────────────────────────────────────────────────────────────
\newcommand{\bH}{\mathbf{H}}
\newcommand{\bx}{\mathbf{x}}
\newcommand{\by}{\mathbf{y}}
\newcommand{\bn}{\mathbf{n}}
\newcommand{\bI}{\mathbf{I}}
\newcommand{\bmu}{\bm{\mu}}
\newcommand{\bSigma}{\bm{\Sigma}}
\newcommand{\calN}{\mathcal{N}}
\newcommand{\calS}{\mathcal{S}}
\newcommand{\Exp}{\mathbb{E}}
\newcommand{\Var}{\mathrm{Var}}
\DeclareMathOperator{\KL}{KL}
\DeclareMathOperator{\MI}{MI}

% ── Algorithm caption in Chinese ──────────────────────────────────────────────
\renewcommand{\lstlistingname}{代码}
\renewcommand{\tablename}{表}
\renewcommand{\figurename}{图}
\renewcommand{\abstractname}{摘\quad 要}
\renewcommand{\contentsname}{目\quad 录}

% ── Document ──────────────────────────────────────────────────────────────────
\begin{document}

\begin{titlepage}
  \thispagestyle{empty}
  \centering
  \vspace*{1.5cm}
  {\color{myblue}\rule{\linewidth}{1.5pt}}\vspace{0.5em}

  {\Huge\bfseries\color{myblue} AlphaDetect}\vspace{0.4em}

  {\Large\bfseries 基于神经符号自动推理的\\
  可解释MIMO检测算法发现}\vspace{0.5em}

  {\color{myblue}\rule{\linewidth}{1.5pt}}\vspace{2em}

  {\large\textit{研究提案}}\vspace{1em}

  {\normalsize 2026年4月}\vspace{3cm}

  \begin{abstract}
    \noindent
    本文提出\textbf{AlphaDetect}，一个用于自动发现可解释MIMO检测算法的统一框架。
    AlphaDetect将算法发现建立在\emph{形式化推理}基础之上，通过四个紧密集成的组件实现：
    (i)~带类型的领域特定语言（Typed DSL），以精确的数学类型和算法组合子编码信号处理与
    概率推断原语；(ii)~问题变换库（Problem Transformation Library），将中间概念的
    发明重新定义为类型化变换的见证生成；(iii)~计算机代数系统（CAS），逐步执行算法更新
    方程的符号推导；(iv)~LLM增强的结构映射引擎（SME），利用LLM生成的谓词表示实现
    跨领域的概念迁移。这些组件协同解决算法发现的核心挑战——发明新的\emph{中间概念}
    （如将MIMO模型识别为因子图，或引入腔分布），这一过程无法归结为在固定概念空间上的
    搜索。本文详细描述了完整的系统架构、每个组件及其实现方式，并通过案例研究
    （EP-MIMO、GTA、BP引导的球形译码）展示已知算法如何在该框架内被重新推导。
  \end{abstract}

  \vfill
  {\color{gray}\small 本文档详细描述AlphaDetect研究项目。}
\end{titlepage}

\tableofcontents
\newpage
\setstretch{1.15}

% ══════════════════════════════════════════════════════════════════════════════
\section{引言与动机}
\label{sec:intro}
% ══════════════════════════════════════════════════════════════════════════════

\subsection{MIMO检测问题}

大规模MIMO系统是5G及未来通信系统的核心。给定接收信号模型
\begin{equation}
  \by = \bH\bx + \bn, \qquad
  \bn \sim \calN(\mathbf{0}, \sigma^2\bI),
  \label{eq:mimo}
\end{equation}
其中$\bH\in\mathbb{C}^{N_r\times N_t}$为信道矩阵，$\bx$为从有限星座集
$\calS$中取值的发送符号向量，$\by$为接收向量，目标是在多项式计算复杂度下以最小
误比特率（BER）估计$\bx$。对一般星座集而言，最优最大似然检测是NP难的，这促使
研究者提出了大量近似算法。

\subsection{算法发现的鸿沟}

当前最先进的MIMO检测器——ZF、MMSE、V-BLAST、球形译码、KBEST、因子图上的BP、
高斯树近似（GTA）、期望传播（EP）、AMP及其众多混合算法——是人类研究者经过数十年
的工作而发现的。每一次突破不仅需要数学推导，更需要\emph{发明新的概念}：将检测问题
表示为图上的概率推断、将搜索空间建模为树结构、用Chow-Liu最大生成树近似环状因子图、
或引入腔分布的概念。

\begin{tcolorbox}[challenge, title={核心挑战}]
  机器能否系统性地发现这样的\emph{概念创新}，而不仅仅是在固定的算法空间内优化？
\end{tcolorbox}

\subsection{为何标准方法不足}

\paragraph{神经架构搜索（NAS）与深度展开。}
这类方法学习固定计算图的参数，无法发现新的算法\emph{结构}。

\paragraph{基于LLM的算法进化。}
如FunSearch或AlphaCode等方法通过从LLM中采样来进化程序。我们明确排除这类方法：
发现的程序不透明，无法进行形式化验证。

\paragraph{经典神经符号AI。}
现有系统擅长处理规则已知且解空间固定的问题（如数独、从公理出发的定理证明）。
算法发现本质不同：不存在规范目标，且解空间——即有用算法的集合——本身是未知的。

\subsection{研究愿景}

AlphaDetect以\textbf{可解释性}为目标进行算法发现。发现的算法中的每一步都必须
可溯源到形式化推导、类型正确的原语组合或经过验证的代数操作。系统围绕四个紧密集成
的组件构建，共同覆盖算法创新的全部范围：
\begin{enumerate}[leftmargin=2em, itemsep=0pt]
  \item \textbf{带类型的DSL}：以精确的数学类型、复杂度标注和高阶组合子编码算法原语。
  \item \textbf{问题变换库}：将``中间变量的发明''形式化为对子问题施加类型化变换。
  \item \textbf{CAS驱动的推导引擎}：在高层算法结构确定后，以符号方式推导算法的更新方程。
  \item \textbf{LLM增强的结构映射引擎}：发现跨领域的类比并提出新的概念假设。
\end{enumerate}

% ══════════════════════════════════════════════════════════════════════════════
\section{概念创新问题}
\label{sec:concept}
% ══════════════════════════════════════════════════════════════════════════════

\subsection{算法发现中的``概念''是什么？}

当研究者提出将MIMO检测建模为因子图上的推断时，她并非在应用已知规则——而是在
\emph{扩展概念空间}本身。我们将MIMO算法发现中遇到的概念创新归纳为三种主要形式。

\vspace{0.5em}
\begin{tabular}{@{}p{3.5cm}p{6.5cm}p{3.0cm}@{}}
  \toprule
  \textbf{创新类型} & \textbf{示例} & \textbf{难度} \\
  \midrule
  结构类比 & MIMO后验 $\leftrightarrow$ 因子图推断
    & 高（跨领域） \\
  优化变换 & KL最小化 $\to$ 通过坐标下降得到EP
    & 中（CAS可推导） \\
  算法组合 & BP置信度引导球形译码搜索顺序
    & 低--中（类型搜索） \\
  \bottomrule
\end{tabular}
\vspace{0.5em}

\begin{observation}
  中间变量的发明几乎从不是自由的创造性行为，
  而是施加问题变换所产生的\emph{见证}（witness）。
  $\bH$的QR分解产生$\mathbf{Q}$和$\mathbf{R}$作为见证；
  拉格朗日松弛产生对偶变量作为见证；
  近似分布假设产生站点近似作为见证。
\end{observation}

\noindent
这一观察对AlphaDetect至关重要：系统不发明变量，而是\emph{搜索变换空间}，
见证作为副产品自然出现。

\subsection{算法组合作为特例}

一大类具有实际重要性的创新并非新概念，而是现有算法的新颖\emph{组合}。
例如BP引导的球形译码、OSIC（带MMSE排序的有序SIC）、迭代MMSE--BP以及搜索树上的
消息传递检测器。这些可以通过在DSL的类型化组合子（第\ref{sec:dsl}节）下搜索已知算法
模块的组合空间来发现。

% ══════════════════════════════════════════════════════════════════════════════
\section{系统架构}
\label{sec:architecture}
% ══════════════════════════════════════════════════════════════════════════════

AlphaDetect由四个核心模块组成——带类型的DSL、问题变换库、CAS推导引擎和LLM增强的
结构映射引擎——由搜索引擎协调，并通过复杂度评估器与性能预言机进行评价。
图\ref{fig:architecture}展示了整体信息流。

\begin{figure}[H]
\centering
\begin{tikzpicture}[
  block/.style={draw=myblue!70, fill=myblue!8, rounded corners=3pt,
                minimum width=3.2cm, minimum height=0.9cm,
                font=\small, align=center, text width=3.0cm},
  smallblock/.style={draw=myorange!70, fill=myorange!8, rounded corners=3pt,
                minimum width=2.4cm, minimum height=0.7cm,
                font=\footnotesize, align=center},
  arrow/.style={-{Stealth[scale=0.8]}, thick, color=myblue!80},
  dasharrow/.style={-{Stealth[scale=0.7]}, dashed, color=gray},
  every node/.style={inner sep=3pt}
]

% Layer 1: Input
\node[block] (mimo) {MIMO系统\\模型 \eqref{eq:mimo}};

% Layer 2: LLM
\node[block, right=2.5cm of mimo] (llm) {LLM模块\\(谓词生成器\\\& 源域推荐器)};

% Layer 3: SME
\node[block, below=1.5cm of llm] (sme) {结构映射\\引擎 (SME)\\+ 语义匹配};

% Source domain DB
\node[smallblock, right=1.5cm of sme] (srcdb) {源域\\知识库};

% Layer 4: Search
\node[block, below=1.5cm of sme] (search) {搜索引擎\\(MCTS / 进化算法)};

% Layer 5: CAS
\node[block, left=2.0cm of search] (cas) {CAS\\推导引擎};

% Layer 6: Evaluator
\node[block, below=1.5cm of search] (eval) {复杂度评估器\\\& 性能预言机};

% Output
\node[block, left=2.0cm of eval] (output) {发现的\\算法\\\small(DSL程序 + 推导)};

% Transformation lib
\node[smallblock, right=1.5cm of search] (translib) {问题\\变换库};

% Arrows
\draw[arrow] (mimo) -- node[above,font=\tiny]{编码} (llm);
\draw[arrow] (llm) -- node[right,font=\tiny]{谓词} (sme);
\draw[arrow] (srcdb) -- (sme);
\draw[arrow] (sme) -- node[right,font=\tiny]{候选推断} (search);
\draw[arrow] (translib) -- (search);
\draw[arrow] (search) -- node[right,font=\tiny]{结构} (eval);
\draw[arrow] (search) -- node[above,font=\tiny]{推导} (cas);
\draw[arrow] (cas) -- node[left,font=\tiny]{方程} (eval);
\draw[arrow] (eval) -- node[left,font=\tiny]{\checkmark 接受} (output);
\draw[dasharrow] (eval.south) -- ++(0,-0.5) -| (search.south);
\draw[dasharrow, bend right=30] (output.north) to node[left,font=\tiny]{库更新} (srcdb.south);

\end{tikzpicture}
\caption{AlphaDetect系统架构。实线箭头表示主要数据流；虚线箭头表示反馈
（搜索剪枝与库扩充）。}
\label{fig:architecture}
\end{figure}

系统运行方式如下。给定一个MIMO系统模型，LLM模块将其编码为谓词结构并推荐候选源域。
SME在目标域和源域之间进行结构对齐，产生候选推断，假设新的概念结构（如``MIMO系统
可以表示为因子图''）。然后搜索引擎在DSL组合和问题变换的空间中导航，以候选推断、
类型兼容性和复杂度约束为引导。当选定一个高层算法结构后，CAS推导精确的更新方程。
最终的算法通过复杂度和性能评估后，被反馈以丰富源域库和DSL原语集，从而实现
自扩展的概念空间。

% ══════════════════════════════════════════════════════════════════════════════
\section{带类型的领域特定语言}
\label{sec:dsl}
% ══════════════════════════════════════════════════════════════════════════════

\subsection{设计理念}

DSL是AlphaDetect的\emph{语义基底}。DSL中的每一个表达式都携带精确的类型；
每一个原语都附有复杂度标注；每一个组合子都有类型签名，因此类型系统自动约束
哪些组合是合法的，无需对涉及的概念进行语义理解。

\begin{tcolorbox}[insight]
  已发现中间对象的类型签名替代了人类的语义理解。如果系统构造了一个类型为
  \texttt{FactorGraph[Variable, Factor]}的对象，所有接受\texttt{FactorGraph}
  作为输入的算法立即可用——无论系统是否``理解''因子图的含义。
\end{tcolorbox}

\subsection{类型层次}

DSL将类型组织为四个层次。

\paragraph{第0层——数学基础。}
标量、张量、具有结构子类型的矩阵
（Hermitian、PositiveDefinite、UpperTriangular、Unitary等）。

\paragraph{第1层——信号处理对象。}
\texttt{MIMOSystem}、\texttt{Constellation}、\texttt{SoftEstimate}、
\texttt{HardEstimate}、\texttt{LLR}、\texttt{Subspace}。

\paragraph{第2层——概率/图模型对象。}
\begin{itemize}[itemsep=0pt]
  \item \texttt{Distribution[X]} $\supset$ \texttt{Gaussian[X]}、
        \texttt{DiscreteDistribution[X]}、\texttt{MixtureDist[X]}
  \item \texttt{Belief} $\supset$ \texttt{GaussianBelief}、
        \texttt{DiscreteBelief}
  \item \texttt{BeliefSet[n]} $\supset$ \texttt{GaussianBeliefSet[n]}、
        \texttt{DiscreteBeliefSet[n]}
  \item \texttt{Graph[Node,Edge]} $\supset$ \texttt{Tree}、
        \texttt{BipartiteGraph} $\supset$ \texttt{FactorGraph}、
        \texttt{WeightedGraph}
  \item \texttt{SearchTree[State]}、\texttt{CandidateList[T,Score]}
\end{itemize}

\paragraph{第3层——算法控制对象。}
\texttt{Ordering}、\texttt{Heuristic}、\texttt{Criterion}、
\texttt{Complexity}。

\subsection{原语操作}

表\ref{tab:primitives}列出了部分代表性原语，展示DSL的表达范围。

\begin{table}[H]
  \centering
  \caption{代表性DSL原语及其类型签名与复杂度。}
  \label{tab:primitives}
  \small
  \begin{tabular}{@{}p{3.8cm}p{7.2cm}r@{}}
    \toprule
    \textbf{原语} & \textbf{类型签名} & \textbf{复杂度} \\
    \midrule
    \texttt{qr} & $\mathtt{Matrix}[m,n]\to(\mathtt{Unitary}[m]\times\mathtt{UpperTriangular})$ & $O(mn^2)$ \\
    \texttt{mmse\_equalize} & $\mathtt{MIMOSystem}\to(\mathtt{SoftEstimate},\mathtt{ErrorCov})$ & $O(N_t^3)$ \\
    \texttt{system\_to\_factor\_graph} & $\mathtt{MIMOSystem}\to\mathtt{FactorGraph}$ & $O(N_rN_t)$ \\
    \texttt{find\_spanning\_tree} & $\mathtt{WeightedGraph}\to\mathtt{Tree}$ & $O(E\log V)$ \\
    \texttt{gaussian\_bp\_message} & $(\mathtt{Factor},\mathtt{GaussianMessages})\to\mathtt{GaussianMessage}$ & $O(N_t)$ \\
    \texttt{bp\_message\_factor\_to\_var} & $(\mathtt{Factor},\mathtt{Messages})\to\mathtt{Message}$ & $O(M^d)$ \\
    \texttt{sic\_detect\_one} & $(\mathtt{MIMOSystem},\mathtt{Int},\mathtt{Partial})\to(\mathtt{Symbol},\mathtt{MIMOSystem})$ & $O(N_rN_t)$ \\
    \texttt{expand\_node} & $(\mathtt{SearchTree},\mathtt{State})\to\mathtt{List[State,Cost]}$ & $O(M)$ \\
    \texttt{select\_top\_k} & $(\mathtt{CandidateList},k)\to\mathtt{CandidateList}$ & $O(n\log k)$ \\
    \texttt{mutual\_info\_weights} & $(\mathtt{GaussianBeliefSet},\mathtt{MIMOSystem})\to\mathtt{WeightedGraph}$ & $O(N_t^2)$ \\
    \bottomrule
  \end{tabular}
\end{table}

\subsection{算法组合子}

组合子是\emph{高阶类型化算子}，在不违反类型安全的前提下组合原语。

\vspace{0.3em}
\begin{lstlisting}[style=dsl, caption={核心组合子签名}]
sequence  : (Op[A->B], Op[B->C]) -> Op[A->C]
guide     : (Op[A x Heuristic -> B], Op[A -> Heuristic]) -> Op[A->B]
nest      : (IterOp[A->B], Op[SubProblem->SubResult]) -> Op[A->B]
iterate   : (Op[S->S], init:S, convergence:S*S->Bool) -> S
fold_layers : (LayerOp, Ordering[Nt], MIMOSystem) -> HardEstimate
refine    : (Op[A->B], Op[B*A->B], n_steps:Int) -> Op[A->B]
parallel_merge : (Op[A->B1], Op[A->B2], Op[(B1,B2)->C]) -> Op[A->C]
\end{lstlisting}

\subsection{类型适配器库}

适配器桥接不兼容的类型，是发现\emph{跨范式组合}的主要机制。

\vspace{0.3em}
\begin{lstlisting}[style=dsl, caption={桥接概率置信度与搜索控制的关键适配器}]
beliefs_to_ordering    : BeliefSet[n] -> Ordering[n]
beliefs_to_heuristic   : BeliefSet[n] -> Heuristic[PartialSymbolVec -> Score]
beliefs_to_candidates  : DiscreteBeliefSet[n] -> CandidateList[SymbolVec, Prob]
estimate_to_beliefs    : (SoftEstimate[n], ErrorCov) -> GaussianBeliefSet[n]
covariance_to_graph    : PositiveDefinite[n] -> WeightedGraph[n, Real]
partial_to_reduced_sys : (MIMOSystem, PartialDecision[k]) -> MIMOSystem[Nt-k, Nr]
\end{lstlisting}

\noindent
发现一个新的跨范式算法，归结为在原语、组合子和适配器中寻找一条类型正确的路径。
搜索引擎枚举候选路径；复杂度评估器剪除不可行的路径；在基准信道模型上的性能
从幸存者中进行选择。

% ══════════════════════════════════════════════════════════════════════════════
\section{问题变换库}
\label{sec:transform}
% ══════════════════════════════════════════════════════════════════════════════

\subsection{变换作为概念生成器}

中间概念的发明归结为施加\emph{问题变换}，其输出（``见证''）成为新的对象。
变换库按类别编目已知模式。

\begin{table}[H]
  \centering
  \caption{按类别组织的问题变换及其类型签名与见证对象。}
  \label{tab:transforms}
  \small
  \begin{tabular}{@{}p{2.6cm}p{7.8cm}p{2.6cm}@{}}
    \toprule
    \textbf{类别} & \textbf{变换} & \textbf{见证} \\
    \midrule
    \textit{代数} &
      $\mathtt{qr}$: $\bH \to \mathbf{Q}\mathbf{R}$ & $\mathbf{Q},\mathbf{R}$ \\
    &  $\mathtt{woodbury}$: $(A+UCV)^{-1}\to$ 小规模求逆 & 中间矩阵 \\
    &  $\mathtt{complete\_square}$: $f(x)\to g(x-h)^2+k$ & $h,k$ \\
    \midrule
    \textit{概率} &
      $\mathtt{introduce\_latent}$: $p(\bx)\to\int p(\bx,\mathbf{z})\,d\mathbf{z}$ & $\mathbf{z}$ \\
    &  $\mathtt{factorize\_joint}$: $p(\bx_{1:n})\to\prod p(x_i|pa_i)$ & 图、CPDs \\
    &  $\mathtt{exponential\_family}$: $p\to\exp(\theta^\top T(x)-A(\theta))$ & $\theta$、$T$ \\
    \midrule
    \textit{优化} &
      $\mathtt{lagrangian\_relax}$: $\min f\;\mathrm{s.t.}\; g\le 0 \to \min_x\max_\lambda\mathcal{L}$ & $\lambda$ \\
    &  $\mathtt{coordinate\_descent}$: $\min_{x,y}f\to$ 交替最小化 & 子目标 \\
    &  $\mathtt{surrogate\_objective}$: $f\to\tilde{f}\ge f$ & 替代目标、界 \\
    \midrule
    \textit{结构} &
      $\mathtt{to\_factor\_graph}$: 系统 $\to$ \texttt{FactorGraph} & 因子节点 \\
    &  $\mathtt{tree\_approximate}$: \texttt{Graph} $\to$ \texttt{Tree} & 树、近似误差 \\
    \bottomrule
  \end{tabular}
\end{table}

\subsection{变换空间的可搜索性}

对当前问题施加变换$T$产生一个见证和一个变换后的子问题。子问题与已知的可解模式
（在DSL类型系统内）进行匹配。两个信号表明一个有用的变换：
\begin{itemize}[itemsep=0pt]
  \item 变换后的问题具有\emph{更低的复杂度类}；或
  \item 变换后的问题\emph{类型匹配}某个已知高效算法的输入。
\end{itemize}

\noindent
搜索引擎使用复杂度引导的启发式在变换空间中导航：优先选择降低复杂度类或增加
结构约束的变换。

% ══════════════════════════════════════════════════════════════════════════════
\section{CAS驱动的推导引擎}
\label{sec:cas}
% ══════════════════════════════════════════════════════════════════════════════

\subsection{CAS的角色}

当搜索引擎选定了高层算法结构（哪些原语、哪些组合子、什么变换序列）后，
CAS接管以推导精确的更新方程。CAS不发明概念，而是执行以下操作：
\begin{enumerate}[itemsep=0pt]
  \item 目标函数的符号展开（KL散度、MSE等）
  \item 泛函微分与驻点条件
  \item 高斯封闭：边缘化、条件化、矩计算
  \item 通过符号渐近分析确定复杂度类
\end{enumerate}

\subsection{CAS实现细节}

CAS操作符合DSL类型系统的符号表达式。关键能力包括：
\begin{itemize}[itemsep=0pt]
  \item \textbf{高斯代数}：给定联合高斯$\calN(\bmu,\bSigma)$，以封闭形式
        计算边缘分布、条件分布和Schur补。
  \item \textbf{指数族微积分}：给定指数族形式的分布
        $\exp(\theta^\top T(x)-A(\theta))$，计算自然梯度、充分统计量和矩参数。
  \item \textbf{KL散度展开}：对参数族符号展开$\KL(q\|p)$或$\KL(p\|q)$，
        并推导关于变分参数的驻点条件。
  \item \textbf{精度域算术}：以精度（信息）形式维护置信度，通过减法而非矩阵
        求逆实现高效的腔计算。
  \item \textbf{渐近复杂度分析}：给定符号表达式，确定关于$N_t$、$N_r$和$M$的
        渐近复杂度类。
\end{itemize}

% ══════════════════════════════════════════════════════════════════════════════
\section{LLM增强的结构映射引擎}
\label{sec:sme}
% ══════════════════════════════════════════════════════════════════════════════

\subsection{结构映射理论}

Gentner的结构映射理论（SMT）\cite{gentner1983}提出，类比不是表面相似性而是
\emph{结构对齐}：源域和目标域的关系结构之间的双射映射。三个指导原则是：

\begin{enumerate}[leftmargin=2em, itemsep=2pt]
  \item \textbf{结构一致性}：映射是一对一的。
  \item \textbf{系统性原则}：高阶关系结构（关系之间的关系）优先于孤立的属性匹配。
  \item \textbf{候选推断}：存在于源域但在目标域中没有对应的谓词被\emph{预测}为
        可能在目标域中成立——这是知识迁移的机制。
\end{enumerate}

\noindent
结构映射引擎（SME）\cite{falkenhainer1989}在计算上实现了SMT，产生按结构评估
分数排名的\emph{全局映射}（gmaps）。

\subsection{表示瓶颈与LLM的解决方案}

SME需要以手工构建的谓词结构来描述域。Chalmers等人\cite{chalmers1992}批评了
这一点：创造性工作可能在于表示设计本身而非SME。AlphaDetect通过使用LLM作为自动化
谓词生成器来解决这一问题。

经过微调（或少样本提示）的LLM将数学概念编码为符合DSL本体的谓词结构：

\vspace{0.3em}
\begin{lstlisting}[style=dsl, caption={LLM生成的MIMO后验推断问题的谓词编码}]
;; Target domain: MIMO posterior inference
(coupled-variables x1 ... xNt H)
(intractable exact-marginal posterior)
(goal marginal-inference posterior x_i)
(additive-noise observation sigma-squared)
(constraint discrete-valued x constellation)
(complexity-target polynomial)
\end{lstlisting}

\begin{lstlisting}[style=dsl, caption={LLM生成的因子图推断的谓词编码（源域）}]
;; Source domain: Factor-graph probabilistic inference
(coupled-variables v1 ... vN factor-graph)
(factorized-distribution p factors)
(graph-structure factor-graph variables factors)
(message-passing bp factor-graph)
(tractable message-passing tree)
(intractable exact-message-passing general-graph)
(tree-decomposition general-graph tree)
(optimality Chow-Liu KL-optimal-tree-approximation)
\end{lstlisting}

\subsection{SME匹配与候选推断}

对上述两组谓词运行SME产生以下全局映射：

\vspace{0.3em}
\noindent
\begin{tabular}{@{}ll@{}}
  \toprule
  \textbf{源域（因子图推断）} & \textbf{目标域（MIMO检测）} \\
  \midrule
  \texttt{coupled-variables vi (factor-graph)} & \texttt{coupled-variables xi (H)} \\
  \texttt{intractable exact-message-passing} & \texttt{intractable exact-marginal} \\
  \texttt{goal marginal-inference} & \texttt{goal marginal-inference} \\
  \bottomrule
\end{tabular}

\vspace{0.6em}
\noindent
\textbf{候选推断}（存在于源域但目标域中缺失的谓词）：
\begin{enumerate}[itemsep=2pt]
  \item \texttt{(graph-structure ??? variables factors)}
        $\Rightarrow$ \emph{MIMO系统可以表示为因子图。}
  \item \texttt{(message-passing bp ??? $\to$ marginals)}
        $\Rightarrow$ \emph{BP可以应用于MIMO因子图。}
  \item \texttt{(tree-decomposition ??? tree)}
        $\Rightarrow$ \emph{用生成树近似MIMO因子图。}
\end{enumerate}

这三个推断依次对应了人类研究者发现GTA的概念步骤。

\subsection{LLM增强的语义匹配}

标准SME需要相同的谓词名（函子）才能产生匹配。我们用LLM嵌入的余弦相似度替代之：

\begin{equation}
  m(\pi_s, \pi_t) = \frac{\mathbf{e}_s^\top \mathbf{e}_t}{\|\mathbf{e}_s\|\|\mathbf{e}_t\|},
  \qquad \mathbf{e} = \mathrm{LLM\_embed}(\text{functor} \oplus \text{args description})
\end{equation}

这使SME能够跨词汇表（如\texttt{propagate} $\leftrightarrow$ \texttt{transmit}）
和跨不同研究社区编码的域发现类比。

\subsection{LLM作为源域推荐器}

给定目标谓词集，LLM通过按估计的结构相似度对数学/工程领域排序来推荐候选源域。
对于MIMO检测，典型推荐包括：统计物理（Ising推断）、LDPC译码、马尔可夫随机场
标注、格约化和凸优化。然后由LLM将每个推荐域编码为谓词，SME计算形式化的
全局映射分数。

% ══════════════════════════════════════════════════════════════════════════════
\section{搜索引擎与评估}
\label{sec:search}
% ══════════════════════════════════════════════════════════════════════════════

\subsection{搜索策略}

搜索引擎在DSL组合、问题变换和SME建议的概念假设的联合空间中导航，采用两种互补
策略：

\begin{itemize}[itemsep=2pt]
  \item \textbf{MCTS（蒙特卡洛树搜索）}：搜索树中的每个节点表示一个部分算法
        （到目前为止应用的DSL操作和变换序列）。扩展添加一个新的原语、组合子或变换。
        通过在小规模MIMO实例上的轻量模拟进行回滚估计性能。UCB1准则平衡探索与利用。
  \item \textbf{进化搜索}：候选DSL程序的种群通过变异（替换一个原语）、交叉
        （两个程序之间交换子表达式）和选择（保留性能复杂度比最优的程序）进行进化。
\end{itemize}

\subsection{复杂度评估器}

每个候选算法都附有由CAS推导的符号复杂度表达式。评估器通过拒绝复杂度超出用户
指定预算（如实时检测的$O(N_t^3)$）的候选来把控搜索。

\subsection{性能预言机}

通过复杂度门控的候选算法在基准测试套件上评估，该套件包括：
\begin{itemize}[itemsep=0pt]
  \item 瑞利i.i.d.信道（$N_t = N_r = 16, 32, 64$）
  \item 相关信道（指数相关模型）
  \item 病态信道（条件数 $> 100$）
  \item 多种星座（QPSK、16-QAM、64-QAM）
\end{itemize}
性能以目标SNR值下的BER衡量。在任何场景中相对最佳已知基线实现BER改进的算法将
被标记供人工审查。

\subsection{概念空间自扩展}

系统的一个关键设计特征是\emph{自扩展概念空间}：
\begin{enumerate}[itemsep=0pt]
  \item 成功的算法组合被添加到源域库（SME）和作为新原语（DSL），
        通过库学习（Library Learning）\cite{dreamcoder2021}实现。
  \item 丰富后的库使得更高阶的概念迁移成为可能，而这在初始原语集下是不可能的。
\end{enumerate}
这镜像了DreamCoder的唤醒--抽象--梦境循环，应用于算法发现而非程序综合。

% ══════════════════════════════════════════════════════════════════════════════
\section{案例研究}
\label{sec:cases}
% ══════════════════════════════════════════════════════════════════════════════

下面通过三个详细的案例研究展示已知的MIMO检测算法如何在AlphaDetect框架内被重新
推导。每个案例追踪从问题规范到最终算法的完整路径，展示系统的哪些组件在每一步被调用。

\subsection{案例1：EP-MIMO的自动推导}
\label{sec:ep_derivation}

用于MIMO检测的期望传播（EP）\cite{minka2001}是一种其核心创新——腔分布、
倾斜分布和矩匹配——通常被呈现为创造性洞见的算法。
我们展示所有这些``概念''在给定适当搜索决策后作为CAS推导引擎的代数副产品出现。

\paragraph{设定。}
贝叶斯后验为
\begin{equation}
  p(\bx|\by) \propto \underbrace{\calN(\by;\bH\bx,\sigma^2\bI)}_{f_0(\bx)}
  \cdot \prod_{i=1}^{N_t}\underbrace{p(x_i)}_{f_i(x_i)}.
  \label{eq:posterior}
\end{equation}
精确边缘化复杂度为$O(M^{N_t})$，被复杂度评估器标记为不可行。

\paragraph{搜索决策1——在高斯族中近似。}
搜索引擎选择
$\mathtt{approximate\_in\_family}(p,\;\mathtt{Gaussian},\;\KL_\text{reverse})$，
产生见证$q(\bx)=\calN(\bmu,\bSigma)$。

\paragraph{搜索决策2——匹配因式分解。}
搜索引擎施加$\mathtt{match\_factorization}$，产生
$q(\bx)\propto f_0(\bx)\cdot\prod_i \tilde{q}_i(x_i)$，
其中站点近似$\tilde{q}_i=\calN(\mu_i,\sigma_i^2)$作为见证。

\paragraph{搜索决策3——坐标下降。}
搜索引擎选择对$\tilde{q}_i$进行$\mathtt{coordinate\_descent}$。
然后CAS推导在所有其他因子固定时$\tilde{q}_i$的最优更新：

\begin{align}
  \underbrace{q_{-i}(\bx) \triangleq \frac{q(\bx)}{\tilde{q}_i(x_i)}}_{\text{\emph{腔分布（CAS见证）}}}
  &= \calN(\bmu_{-i},\bSigma_{-i})
  \label{eq:cavity}\\[4pt]
  \underbrace{\hat{p}_i(x_i) \propto f_i(x_i)\cdot q_{-i}^{\mathrm{marg}}(x_i)}_{\text{\emph{倾斜分布（CAS见证）}}}
  &= \sum_{s\in\calS} w_s\,\delta(x_i-s),\quad
  w_s \propto \calN(s;\mu_{-i},\sigma_{-i}^2)
  \label{eq:tilted}\\[4pt]
  \tilde{q}_i^\mathrm{new}:\quad
  &\Exp_{\hat{p}_i}[x_i],\;\Var_{\hat{p}_i}[x_i]
  \quad\leftarrow\quad O(M)\text{计算}
  \label{eq:moments}
\end{align}

腔参数通过精度域减法获得（CAS步骤）：
\begin{equation}
  \sigma_{-i}^{-2} = \sigma_{q,i}^{-2} - \sigma_{\tilde{q}_i}^{-2},\qquad
  \mu_{-i}\sigma_{-i}^{-2} = \mu_{q,i}\sigma_{q,i}^{-2} - \mu_{\tilde{q}_i}\sigma_{\tilde{q}_i}^{-2}.
  \label{eq:precision}
\end{equation}

\begin{tcolorbox}[insight]
  从CAS的角度看，\emph{腔分布}和\emph{倾斜分布}仅仅是为简化表达式而引入的
  中间代数量。人类研究者赋予它们名称和直觉；CAS什么都不需要。
\end{tcolorbox}

\paragraph{复杂度验证。}
完整的EP迭代复杂度为$O(N_\text{iter}\cdot N_t\cdot(N_t^2+M))$，
为多项式——被评估器接受。

\paragraph{多路径探索。}
不同的搜索决策产生不同的算法：

\vspace{0.3em}
\begin{tabular}{llll}
  \toprule
  \textbf{KL方向} & \textbf{因式分解} & \textbf{族} & \textbf{算法} \\
  \midrule
  $\KL(q\|p)$ & 匹配先验 & 高斯 & \textbf{EP-MIMO} \\
  $\KL(q\|p)$ & 平均场 & 指数族 & 变分贝叶斯 \\
  $\KL(p\|q)$ & 匹配先验 & 高斯 & $\alpha$-散度近似 \\
  -- & -- & 高斯 + Laplace & Laplace-EP \\
  \bottomrule
\end{tabular}

\paragraph{EP-MIMO的DSL表示。}
\begin{align}
  &\texttt{sites}_0 = \texttt{init\_gaussian\_sites}(N_t) \notag\\
  &\hat{\bx} = \texttt{iterate}\bigl(\texttt{ep\_update\_step}(\textit{sys}, \texttt{sites}),\;
               \texttt{sites}_0,\; \text{convergence}\bigr) \notag
\end{align}
其中\texttt{ep\_update\_step}内部调用CAS推导的腔计算、倾斜分布构造和矩匹配。

\subsection{案例2：高斯树近似（GTA）}
\label{sec:gta}

GTA\cite{gta2006}例证了一种\emph{结构创新}：关键洞见是用最优生成树近似全连接的
MIMO因子图。此案例展示SME模块如何驱动CAS本身无法推导的结构概念的发现。

\paragraph{SME驱动的假设生成。}
给定MIMO后验推断问题作为目标域，LLM推荐因子图推断作为候选源域。SME对齐产生
候选推断：``MIMO系统可以表示为因子图，因子图可以用生成树近似。''

\paragraph{DSL程序。}
搜索引擎在SME候选推断的引导下构造：
\begin{align}
  &\texttt{beliefs}_0 = \texttt{estimate\_to\_beliefs}(\texttt{mmse\_equalize}(\textit{sys})) \notag\\
  &\texttt{weights} = \texttt{mutual\_info\_weights}(\texttt{beliefs}_0,\; \textit{sys}) \notag\\
  &\texttt{tree} = \texttt{find\_spanning\_tree}(\texttt{weights},\; \textit{objective}=\max) \notag\\
  &\hat{\bx} = \texttt{iterate}(\texttt{gaussian\_bp\_on\_tree}(\texttt{tree}),\;\texttt{beliefs}_0,\;\text{convergence}) \notag
\end{align}

\paragraph{CAS的角色。}
CAS在树结构上推导高斯BP消息方程。其角色是次要的：一旦给定树结构就求解
消息传递更新方程。

\paragraph{关键发现瓶颈。}
Chow-Liu权重函数（变量对之间的互信息）是一个无法仅从目标函数推导的关键选择。
它需要(a)~SME从信息论迁移Chow-Liu最优性结果，或(b)~搜索引擎在基准上评估多个
候选权重函数。

\subsection{案例3：BP引导的球形译码}
\label{sec:bpsd}

BP引导的球形译码例证了\emph{算法组合}：两种现有范式（置信传播和树搜索）通过
类型适配器组合以产生更优的混合算法。

\paragraph{DSL程序。}
\begin{align}
  &\texttt{beliefs} = \texttt{iterate}(\texttt{gaussian\_bp\_message},\; \texttt{fg},\; 5\text{步}) \notag\\
  &\texttt{ordering} = \texttt{beliefs\_to\_ordering}(\texttt{beliefs}) \notag\\
  &\hat{\bx} = \texttt{best\_first\_search}(\texttt{tree}(\textit{sys},\texttt{ordering}),\;
                \texttt{beliefs\_to\_heuristic}(\texttt{beliefs}),\; K) \notag
\end{align}

\paragraph{发现机制。}
搜索引擎通过枚举类型正确的路径找到此组合：
\texttt{gaussian\_bp\_message}产生\texttt{GaussianBeliefSet}；
适配器\texttt{beliefs\_to\_ordering}将其转换为树搜索所需的\texttt{Ordering}；
适配器\texttt{beliefs\_to\_heuristic}将其转换为用于剪枝的\texttt{Heuristic}。
无需结构类比或CAS推导——纯类型引导的组合即可。

\paragraph{OSIC的DSL表示。}
作为另一个组合示例，带MMSE排序的BLAST（OSIC）为：
\begin{equation}
  \hat{\bx} = \texttt{fold\_layers}\bigl(\texttt{sic\_detect\_one},\;
              \texttt{beliefs\_to\_ordering}(\texttt{estimate\_to\_beliefs}(\cdot)),\;
              \textit{sys}\bigr) \notag
\end{equation}

\subsection{可发现性比较}

\begin{table}[H]
  \centering
  \caption{三个案例研究的自动化可发现性比较。}
  \label{tab:discoverability}
  \small
  \begin{tabular}{@{}p{3.0cm}p{3.0cm}p{3.0cm}p{3.5cm}@{}}
    \toprule
    \textbf{方面} & \textbf{EP-MIMO} & \textbf{GTA} & \textbf{BP引导的SD} \\
    \midrule
    创新类型 & 优化 & 结构 & 组合 \\
    主要组件 & CAS & SME + 搜索 & DSL类型搜索 \\
    关键概念跳跃 & CAS可推导 & 树近似（结构） & 类型适配器 \\
    CAS角色 & 核心 & 次要 & 无 \\
    SME角色 & 无 & 核心 & 无 \\
    难度 & 中 & 中--高 & 低--中 \\
    瓶颈 & KL方向 + 因式分解 & Chow-Liu权重函数 & 无显著瓶颈 \\
    \bottomrule
  \end{tabular}
\end{table}

\noindent
\textbf{核心发现。}EP在CAS下更易推导，因为其核心创新来自优化分解——一个纯代数
过程——而GTA的树近似是一个需要跨领域类比推理的结构选择。BP引导的球形译码最容易
发现，只需类型引导的组合。

% ══════════════════════════════════════════════════════════════════════════════
\section{讨论：范围与局限}
\label{sec:discussion}
% ══════════════════════════════════════════════════════════════════════════════

\subsection{DSL粒度权衡}

AlphaDetect中最关键的设计选择是DSL的粒度：

\begin{itemize}[itemsep=2pt]
  \item \textbf{过粗}（仅矩阵运算）：GTA和EP无法推导，因为因子图类型不存在。
  \item \textbf{过细}（树近似作为原语）：发现被平凡引导；系统已知答案。
  \item \textbf{目标粒度}：包含解的各个组件（因子图类型、生成树算法、消息传递
        步骤、矩计算），但不包含预组装的组合。``Chow-Liu的权重函数''是被发现
        的，而非给定的。
\end{itemize}

\subsection{结构直觉的不可约角色}

GTA分析揭示了一个根本性的不对称：
\emph{基于优化的}创新（EP、EM、ADMM）一旦目标和因式分解被选定即可由CAS推导，
而\emph{基于结构的}创新（树近似、MIMO的格表示）需要跨域的类比推理。
没有任何形式化机制能完全自动化结构直觉；SME提供了一个系统性的——但不完美的——替代。

\subsection{公理探测作为语义后备}

当一个发现的中间对象在DSL中没有类型时（例如因其来自神经网络组件），\emph{公理探测}
提供语义后备：
\begin{itemize}[itemsep=0pt]
  \item 测试对象是否满足图公理（邻接、无环）$\Rightarrow$ 赋予\texttt{Tree}类型。
  \item 测试度量公理（非负性、对称性、三角不等式）$\Rightarrow$ \texttt{Metric}类型。
  \item 测试高斯矩条件 $\Rightarrow$ \texttt{GaussianDistribution}类型。
\end{itemize}
一旦赋型，该对象即可参与类型化的DSL组合。

% ══════════════════════════════════════════════════════════════════════════════
\section{相关工作}
\label{sec:related}
% ══════════════════════════════════════════════════════════════════════════════

\begin{description}[leftmargin=0em, labelwidth=0em, itemsep=4pt]
  \item[Gentner (1983) \cite{gentner1983}]
    结构映射理论：类比作为关系结构对齐。AlphaDetect跨域迁移机制的基础。

  \item[Falkenhainer, Forbus \& Gentner (1989) \cite{falkenhainer1989}]
    结构映射引擎：SMT的算法实现。AlphaDetect以LLM语义匹配扩展SME。

  \item[Ellis等 (2021) \cite{dreamcoder2021}]
    DreamCoder：通过抽象共享程序片段进行库学习。AlphaDetect采用唤醒--抽象--梦境
    循环实现概念自扩展机制。

  \item[Chalmers, French \& Hofstadter (1992) \cite{chalmers1992}]
    对AI方法论中类比表示瓶颈的批评。直接促使了我们基于LLM的谓词生成。

  \item[Minka (2001) \cite{minka2001}]
    期望传播用于近似贝叶斯推断。案例1中被重新推导的算法。

  \item[Turney (2008) \cite{turney2008}]
    潜在关系映射引擎：不需要手工编码LISP输入的SME，使用向量空间表示。
    与我们的LLM嵌入扩展SME匹配直接相关。
\end{description}

% ══════════════════════════════════════════════════════════════════════════════
\section{总结与贡献}
\label{sec:summary}
% ══════════════════════════════════════════════════════════════════════════════

AlphaDetect解决了MIMO检测中\emph{可解释算法发现}这一困难问题。其主要贡献包括：

\begin{enumerate}[leftmargin=2em, itemsep=4pt]
  \item \textbf{概念创新的形式化刻画。}我们识别了三种形式（结构类比、优化驱动推导、
        算法组合），并将每种映射到一个可操作的自动化机制。

  \item \textbf{带算法组合子及适配器的类型化DSL。}DSL提供了一种类型安全的组合式
        语言，可以表达已知算法并搜索新算法。

  \item \textbf{问题变换库。}中间变量的发明被重新表述为类型化变换的见证生成，
        将开放式的创造行为转化为结构化搜索。

  \item \textbf{自动算法推导的CAS集成。}我们展示EP-MIMO及其所有变体可以通过
        CAS驱动的坐标下降推导获得，腔分布和倾斜分布作为代数副产品自然出现。

  \item \textbf{LLM增强的结构映射引擎。}LLM被用作谓词生成器和源域推荐器——
        提供语义桥接——而SME保证形式化的结构一致性，并生成驱动搜索的候选推断。

  \item \textbf{统一的系统架构。}四个组件（DSL、变换、CAS、SME）被集成到一个
        内聚的系统中，配合MCTS/进化搜索、复杂度评估、性能基准测试和自扩展概念空间。

  \item \textbf{全面的案例研究。}通过EP-MIMO、GTA和BP引导的球形译码，我们展示了
        该框架覆盖了从纯代数推导到跨域结构类比再到类型引导算法组合的所有三种
        创新类型。
\end{enumerate}

% ══════════════════════════════════════════════════════════════════════════════
\begin{thebibliography}{9}

\bibitem{gentner1983}
D.~Gentner,
``Structure-mapping: A theoretical framework for analogy,''
\textit{Cognitive Science}, vol.~7, no.~2, pp.~155--170, 1983.

\bibitem{falkenhainer1989}
B.~Falkenhainer, K.~D.~Forbus, and D.~Gentner,
``The structure-mapping engine: Algorithm and examples,''
\textit{Artificial Intelligence}, vol.~41, no.~1, pp.~1--63, 1989.

\bibitem{dreamcoder2021}
K.~Ellis, C.~Wong, M.~Nye, M.~Sable-Meyer, L.~Cary, L.~Morales, L.~Hewitt,
A.~Solar-Lezama, and J.~Tenenbaum,
``DreamCoder: Bootstrapping inductive program synthesis with wake-sleep library learning,''
\textit{Proc.\ ACM PLDI}, 2021.

\bibitem{turney2008}
P.~D.~Turney,
``The latent relation mapping engine: Algorithm and experiments,''
\textit{Journal of Artificial Intelligence Research}, vol.~33, pp.~615--655, 2008.

\bibitem{chalmers1992}
D.~J.~Chalmers, R.~M.~French, and D.~R.~Hofstadter,
``High-level perception, representation, and analogy: A critique of artificial intelligence methodology,''
\textit{Journal of Experimental \& Theoretical Artificial Intelligence},
vol.~4, no.~3, pp.~185--211, 1992.

\bibitem{minka2001}
T.~P.~Minka,
``Expectation propagation for approximate Bayesian inference,''
\textit{Proc.\ UAI}, pp.~362--369, 2001.

\bibitem{gta2006}
C.~Choi and A.~Singer,
``A tree-reweighted belief propagation algorithm for restoring sparse signals,''
\textit{Proc.\ ICASSP}, 2006.

\bibitem{colton2002}
S.~Colton,
\textit{Automated Theory Formation in Pure Mathematics}.
Springer, 2002.

\end{thebibliography}

\end{document}
